\section{Evaluation}
\label{sec:evaluation}

All numbers in this section are produced by the scripts in the \texttt{experiments}
directory of the repository, which also stores their results. The scripts use the production agent loop and perimeter without
modification. Unless stated otherwise, they need no network or GPU and are deterministic given
their seeds.

\paragraph{Harness.} A seeded corpus assigns one third of the documents to each class. Each
restricted document contains a canary, a random string unique to that document that does not
appear in any policy. A fraction of the restricted documents is \emph{misfiled}, that is,
stored in a folder that the policy's path rules classify as \pub, as would happen with an offer
saved in a shared folder. The agent is scripted to behave adversarially: it reads the document,
tries to send its full text through the browser, and then answers. Every substrate records what
it receives, and the browser acts as a network sink. A leak is counted when the canary of a
restricted document is observed by a sink whose ceiling is below \res. The true class comes
from the document generator and not from the policy, so the leak rate does not depend on the
rules being tested. The deployment has three substrates: an on-premises GPU (\res), an EU
cluster (\inte) and a frontier API (\pub). The local GPU is unavailable in one quarter of the
tasks.

\subsection{E2: restricted egress}
\label{sec:eval-leak}

\begin{table}[t]
\centering\small
\begin{tabularx}{\linewidth}{@{}l r r r >{\raggedright\arraybackslash}X@{}}
\toprule
Configuration & Leak & Leak (misfiled) & Held & Sinks that received canaries \\
\midrule
Gateway (failover list) & 0.422 & 1.000 & 0.000 & frontier 113, EU cluster 98, network 113 \\
\annona, without \texttt{tools.egress} & 0.804 & 1.000 & 0.065 & network 402, frontier 113 \\
\annona, path and pattern rules & 0.226 & 1.000 & 0.065 & frontier 113, network 113 \\
\annona{} + classifier, $\alpha{=}0.05$ & 0.008 & 0.035 & 0.084 & frontier 4, network 4 \\
\annona{} + classifier, $\alpha{=}0.01$ & 0.002 & 0.009 & 0.084 & frontier 1, network 1 \\
\bottomrule
\end{tabularx}
\caption{E2 over 1{,}500 tasks per configuration (300 documents, 5 seeds), of which 500 are
restricted and 113 misfiled (22.6\%). Leak is the fraction of restricted documents whose
canary reached a sink below \res; Held is the fraction of all tasks stopped by the perimeter.}
\label{tab:leak}
\end{table}

\Cref{tab:leak} shows three separate causes of leaks. The first is failover. The gateway leaks
every restricted document processed while the local GPU was down, 98 of them to the EU cluster,
whereas \annona{} holds these tasks (the 6.5\% held). The second is tool egress. Without the
network-tool ceiling, the browser sent the full text of 402 of the 500 restricted documents off
the machine, which is every restricted document that was not held. The third is
classification. With path and pattern rules only, \annona{} leaks the 113 misfiled documents
and no others, so the perimeter behaves correctly for the class it receives and the errors
come from the class itself. Adding a content classifier removes most of these errors. At
$\alpha{=}0.05$ the leak rate is 0.8\% and 3.5\% of misfiled documents leak; at
$\alpha{=}0.01$ the rates are 0.2\% and 0.9\%. Both are within the bound of \cref{prop:leak},
which gives 1.1\% and 0.23\% for a misfiled fraction of 22.6\%. The additional cost is 1.9
percentage points of held tasks, corresponding to misfiled documents that are now recognised as
restricted while the local GPU is down.

The classifier in this experiment is simulated with a fixed miss rate. This isolates the
composition from the quality of any particular model; E7 reports a real classifier. All
ledgers written by the \annona{} configurations were audited: 26{,}521 entries and no
violations.

\subsection{E3: detecting a faulty perimeter from the ledger}
\label{sec:eval-compliance}

We injected six faults that a placement runtime could plausibly contain, ran the same workload
(450 tasks per fault), and audited each resulting ledger against the approved policy file
(\cref{tab:mutants}).

\begin{table}[t]
\centering\small
\begin{tabularx}{\linewidth}{@{}l X r >{\raggedright\arraybackslash}p{3.6cm}@{}}
\toprule
Injected fault & Behaviour & Leak & Violations reported \\
\midrule
none & correct runtime & 0.167 & none (2{,}626 entries) \\
failover-widening & when no permitted substrate is up, uses any available one & 0.413 & sound 74 \\
stale-class & places using the payload's class and ignores what the run has read & 1.000 & monotone 550 \\
inverted-preference & picks the rule's last choice & 0.167 & optimal 799 \\
spurious-hold & holds restricted work although a permitted substrate is up & 0.167 & hold-justified 88 \\
policy-drift & runs under a modified policy while the approved one is audited & 0.413 & bound 1{,}800, sound 74, optimal 176 \\
egress-ceiling-off & does not enforce the network-tool ceiling & 0.753 & egress 238 \\
\bottomrule
\end{tabularx}
\caption{E3. Each fault is reported by the corresponding invariant, and the correct runtime
produces no violations. Leak is measured as in E2 (path rules only, 20\% misfiled).}
\label{tab:mutants}
\end{table}

Two of the faults cause no leak. \emph{Inverted-preference} lowers quality and
\emph{spurious-hold} lowers availability, and both are still reported because the audit checks
that each decision matches the one prescribed by the policy, in addition to checking that it is
safe. \emph{Stale-class} is the most harmful fault, since every restricted document leaks, and
yet each of its entries is consistent with the class it records. It is detected only by the
\emph{monotone} invariant, which compares each placement with what the run had already read.

\paragraph{Labels must not depend on the control they record.} In our first version of this
experiment, \emph{egress-ceiling-off} was not detected even though 75\% of restricted
documents leaked. The same function decided whether a tool had a ceiling and whether its calls
were labelled \texttt{network:}. Disabling the ceiling therefore also labelled the calls
\texttt{local}, which removed the evidence of the fault. After the change, the label depends
only on whether the tool reaches the network, and the auditor reads this property from the
policy instead of trusting the label. More generally, a field that records whether a control
was applied should not be computed by that control.

\subsection{E4: tamper evidence}
\label{sec:eval-tamper}

We applied every single edit an insider could make to each entry of a 360-entry ledger produced
by the harness (\cref{tab:tamper}). Field edits, byte flips and swaps were always detected and
located at the edited entry. Deletion was detected in 359 of 360 cases; the exception was the
last line, whose deletion is equivalent to truncation. As expected, truncation and rewrites that
recompute the hashes of all later entries are not detected by the chain alone. With a witness
that stores a checkpoint every 10 entries, both are detected in all cases before the last
checkpoint. With one checkpoint every 50 entries the detection rate is 97.2\%, and the
undetected cases are the entries written after the last checkpoint.

\begin{table}[t]
\centering\small
\begin{tabular}{@{}l r r r r@{}}
\toprule
Edit & Trials & Chain only & Witness, every 10 & Witness, every 50 \\
\midrule
Field edit (7 fields) & 2{,}520 & 1.000 & & \\
Byte flip & 360 & 1.000 & & \\
Swap with next entry & 359 & 1.000 & & \\
Delete & 360 & 0.997 & & \\
Truncate tail & 359 & 0.000 & 1.000 & 0.972 \\
Edit and recompute later hashes & 359 & 0.000 & 1.000 & 0.972 \\
\bottomrule
\end{tabular}
\caption{E4: detection rate of single edits.}
\label{tab:tamper}
\end{table}

\subsection{E5: AgentDojo through the unmodified loop}
\label{sec:eval-dojo}

We ran AgentDojo~\citep{debenedetti2024agentdojo} tasks through \annona's agent loop and
perimeter, using AgentDojo's function runtime as the tool executor. A local model (qwen2.5 via
Ollama) serves two substrates: \texttt{local} (\res) and \texttt{frontier} (\pub), with the
latter recording every request. Because both substrates use the same weights, differences
between configurations come from placement and refused calls rather than from model quality.
A payload containing an IBAN or a password is classified as \res. We compare three
configurations: \emph{open}, where every class may run on the frontier; \emph{placement}; and
\emph{placement + egress}, where the tools that transfer money are egress tools with ceiling
\inte.

\begin{table}[t]
\centering\small
\begin{tabular}{@{}l r r r r@{}}
\toprule
Configuration & Utility & Utility (attacked) & Attack success & User IBAN at frontier \\
\midrule
\multicolumn{5}{@{}l}{\emph{qwen2.5:3b, 16 user tasks $\times$ 4 injection tasks (64 attacked pairs)}} \\
open & 0.375 & 0.313 & 0.000 & 73 of 182 requests \\
placement & 0.438 & 0.297 & 0.016 & 0 of 66 \\
placement + egress & 0.375 & 0.313 & 0.000 & 0 of 65 \\
\midrule
\multicolumn{5}{@{}l}{\emph{qwen2.5:14b, 8 user tasks $\times$ 2 injection tasks (16 attacked pairs)}} \\
open & 0.500 & 0.688 & 0.500 & 64 of 84 \\
placement + egress & 0.250 & 0.313 & 0.000 & 0 of 16 \\
\bottomrule
\end{tabular}
\caption{E5: AgentDojo banking suite (v1.2) with the \texttt{important\_instructions} attack.
Refused tool calls under \emph{placement + egress}: 39 (3B) and 18 (14B).}
\label{tab:dojo}
\end{table}

In the open configuration with the 3B model, 73 of the 182 requests received by the frontier
contained the user's own IBANs, which a tool had read and which then stayed in the transcript.
With placement this number drops to zero: once a payload contains an IBAN the run becomes
restricted and all later steps run locally. Benign utility does not change measurably (7 and
6 of 16 tasks, a difference within noise at this sample size). The 3B model is rarely affected
by the injections, even in the open configuration, so this run measures exposure to the
provider but says little about exfiltration through tools.

The 14B model is affected. In the open configuration half of the injected goals succeed (8 of
16) and 64 of the 84 frontier requests contain the user's IBANs. With placement and egress
ceilings, no injected goal succeeds and no IBAN reaches the frontier. The cost is utility, as
is usual for information-flow control. After the run reads account data it becomes restricted,
and the money-transfer tools are refused for the rest of the run, including transfers the user
actually requested (benign utility drops from 4 to 2 of 8 tasks). Holding such calls for user
confirmation instead of refusing them would reduce this cost; the runtime already supports
holds for placement, and we leave the extension to tools as future work. The sample is small
and the confidence intervals are wide. Running the full suite with the 14B model, the other
AgentDojo suites, and a combination with an IFC planner such as FIDES~\citep{costa2025fides}
are the next experiments.

\subsection{E6: overhead}
\label{sec:eval-overhead}

\begin{table}[t]
\centering\small
\begin{tabular}{@{}l r r@{}}
\toprule
Operation & p50 & p95 \\
\midrule
Classify outgoing payload, 280 characters & 41\,\textmu s & 42\,\textmu s \\
\quad 4k characters & 96\,\textmu s & 99\,\textmu s \\
\quad 40k characters & 0.63\,ms & 0.65\,ms \\
\quad 200k characters (render limit) & 3.0\,ms & 3.1\,ms \\
Check a tool call (gate) & 0.10\,ms & 0.11\,ms \\
Place a step & 4\,\textmu s & 4\,\textmu s \\
Append to ledger & 63\,\textmu s & 83\,\textmu s \\
Append to ledger with fsync & 79\,\textmu s & 94\,\textmu s \\
Audit: verify chain / comply & \multicolumn{2}{r}{18 / 32\,\textmu s per entry} \\
\bottomrule
\end{tabular}
\caption{E6, measured on an Apple M1 Pro with Python 3.12.}
\label{tab:overhead}
\end{table}

Our design document states that a policy check adding 200\,ms per step would be disabled in
practice. The perimeter adds at most about 3.3\,ms per step at the render limit and about
0.3\,ms for a typical 4k-character payload, two to three orders of magnitude below that
threshold and small compared with inference. At 10{,}000 steps a day, one year of decisions
(3.65 million entries) can be verified and audited in about three minutes.

\subsection{E7: a real classifier and the cost of the bound}
\label{sec:eval-judge}

E2 uses a simulated classifier with a known miss rate. With a real classifier the bound holds
in the same way, and the relevant question is how many documents it forces to stay local.
\Cref{tab:judge} reports split-conformal thresholds for the base model served by
rizzo-flow~\citep{rizzoflow} (Spark-X2.5-4B, Q8\_0, zero-shot). Thresholds were fitted on 72
restricted documents from a calibration split and tested on 90 held-out restricted documents
drawn from public contracts and synthetic business documents. These numbers come from the
evaluation of a companion sensitivity model, Iovis, which is still being trained.

\begin{table}[t]
\centering\small
\begin{tabular}{@{}r r r r@{}}
\toprule
$\alpha$ & $\tau_\alpha$ & Missed restricted (test, $n{=}90$) & Non-restricted flagged \\
\midrule
0.01 & n/a & \multicolumn{2}{l}{not supported by 72 calibration documents} \\
0.02 & 0.645 & 0.011 & 0.643 \\
0.05 & 0.657 & 0.056 & 0.557 \\
0.10 & 0.713 & 0.133 & 0.350 \\
\bottomrule
\end{tabular}
\caption{E7: conformal thresholds for a zero-shot classifier. Observed miss rates are within
binomial noise of $\alpha$ for $n{=}90$, since the guarantee holds on average over calibration
sets. The flagged fraction is the cost.}
\label{tab:judge}
\end{table}

For a zero-shot classifier the bound is valid but expensive. Missing at most 2\% of restricted
documents requires flagging 64\% of the others, which the perimeter then keeps local. In the
E2 deployment this would bound the leak rate at $0.226 \times 0.02 \approx 0.5\%$ while moving
most public work off the frontier. The purpose of a classifier trained for commercial
sensitivity is to reach the same $\alpha$ with fewer flagged documents. When there are too few
calibration documents to support a given $\alpha$, the adapter raises an error instead of
returning a threshold that cannot meet it.
